%% Section 2: Classical Littlewood-Richardson Coefficients
%% To be \input{} into the main document.

\section{Classical Littlewood-Richardson Coefficients}
\label{sec:classical}

The Littlewood-Richardson coefficients $c^{\nu}_{\lambda,\mu}$ are
nonnegative integers that appear, independently and without apparent
conspiracy, in three different branches of mathematics.  That these
three appearances yield the \emph{same} number is already a theorem
worth contemplating.  We begin with the three definitions, then survey
the combinatorial models that make the coefficients computable, and
finally recast the entire story in the language of Hopf algebras and
transfer operators --- the doorway to integrability.

Throughout this section, our running example will be the coefficient
\[
  c^{(3,2,1)}_{(2,1),\,(2,1)} = 2.
\]
This is the smallest instance where the coefficient exceeds~$1$, and it
will accompany us into the lattice model of Section~3.

\subsection{Definitions and context}
\label{sec:classical-defn}

A \emph{partition} $\lambda = (\lambda_1 \ge \lambda_2 \ge \cdots \ge
\lambda_k > 0)$ is a weakly decreasing sequence of positive integers.
Its \emph{Young diagram} is the array of left-justified boxes with
$\lambda_i$ boxes in row~$i$; one may picture it as a skyline read
from left to right, each row a terrace one step narrower than the
last.  The \emph{size} $|\lambda| = \sum_i \lambda_i$ is the total
number of boxes.

\begin{definition}[LR coefficients: symmetric functions]
\label{def:lr-symm}
Let $\Lambda$ denote the ring of symmetric functions over $\mathbb{Z}$,
with its basis of Schur functions $\{s_\lambda\}$ indexed by partitions.
The \emph{Littlewood-Richardson coefficients} $c^{\nu}_{\lambda,\mu}$
are the structure constants of this basis:
\[
  s_\lambda \cdot s_\mu
  = \sum_{\nu} c^{\nu}_{\lambda,\mu}\, s_\nu.
\]
Since the Schur functions form a $\mathbb{Z}$-basis of $\Lambda$, the
coefficients are uniquely determined.  That they are
\emph{nonneg\-ative} is a deep fact, not at all obvious from this
definition.
\end{definition}

\begin{definition}[LR coefficients: representation theory]
\label{def:lr-rep}
Let $V_\lambda$ denote the irreducible polynomial representation of
$GL_n(\mathbb{C})$ of highest weight $\lambda$ (equivalently, the
Schur module associated to $\lambda$).  Then
\[
  V_\lambda \otimes V_\mu
  \cong \bigoplus_{\nu} V_\nu^{\oplus\, c^{\nu}_{\lambda,\mu}}
\]
as $GL_n$-representations.  Here the nonnegativity of the
coefficients is manifest: $c^{\nu}_{\lambda,\mu}$ is the multiplicity
of the irreducible $V_\nu$ in a tensor product, hence a dimension,
hence a nonneg\-ative integer.
\end{definition}

\begin{definition}[LR coefficients: Schubert calculus]
\label{def:lr-schub}
Let $\mathrm{Gr}(k,n)$ denote the Grassmannian of $k$-planes in
$\mathbb{C}^n$, and let $\sigma_\lambda \in H^*(\mathrm{Gr}(k,n))$
denote the Schubert class associated to a partition $\lambda$ fitting
inside a $k \times (n-k)$ box.  Then
\[
  \sigma_\lambda \cdot \sigma_\mu
  = \sum_\nu c^{\nu}_{\lambda,\mu}\, \sigma_\nu
  \quad \text{in } H^*(\mathrm{Gr}(k,n)).
\]
The geometric content: $c^{\nu}_{\lambda,\mu}$ counts the number of
points in a triple intersection of Schubert varieties in general
position.  Nonnegativity is again manifest --- one is counting
points.
\end{definition}

\begin{remark}
The fact that Definitions~\ref{def:lr-symm},~\ref{def:lr-rep},
and~\ref{def:lr-schub} give the same numbers is a consequence of the
Borel-Weil theorem and the classical relationship between the
representation ring of $GL_n$ and the cohomology of Grassmannians
(see \cite{fulton-1997} for a unified treatment).  But the
\emph{feeling} of the three definitions is quite different: in
$\Lambda$ the coefficients are algebraic, in representation theory
they are multiplicities, and in Schubert calculus they are
enumerative-geometric.  That a single integer can simultaneously be
all three things is what makes the LR coefficient a meeting place for
different parts of mathematics.
\end{remark}

The coefficients enjoy several structural properties that constrain
and enrich their study:

\begin{itemize}
\item \textbf{Symmetry.}
$c^{\nu}_{\lambda,\mu} = c^{\nu}_{\mu,\lambda}$, since multiplication
in $\Lambda$ is commutative.

\item \textbf{Vanishing.}
$c^{\nu}_{\lambda,\mu} = 0$ unless $|\nu| = |\lambda| + |\mu|$ and
$\nu \supseteq \lambda$, $\nu \supseteq \mu$ (containment of Young
diagrams).

\item \textbf{Saturation} (Knutson-Tao \cite{knutson-tao-1999}).
$c^{N\nu}_{N\lambda,\, N\mu} > 0$ if and only if
$c^{\nu}_{\lambda,\mu} > 0$.  In other words, the \emph{support} of
the LR coefficients is a saturated cone.  This resolved the
longstanding saturation conjecture and plays a key role in the Horn
inequalities characterizing eigenvalues of sums of Hermitian matrices.
\end{itemize}


\subsection{Combinatorial models}
\label{sec:classical-models}

The nonnegativity of $c^{\nu}_{\lambda,\mu}$ cries out for a
combinatorial proof: exhibit a set whose cardinality is
$c^{\nu}_{\lambda,\mu}$.  In fact, several such sets are known, each
illuminating a different aspect of the coefficient.

\subsubsection*{Skew tableaux and the Littlewood-Richardson rule}

Given partitions $\lambda \subseteq \nu$ (meaning $\lambda_i \le
\nu_i$ for all $i$), the \emph{skew shape} $\nu/\lambda$ is the set
of boxes in $\nu$ but not in $\lambda$ --- a diagram with a notch cut
from its northwest corner.

\begin{definition}[Semistandard Young tableau]
A \emph{semistandard Young tableau} (SSYT) of shape $\nu/\lambda$ and
\emph{content} $\mu = (\mu_1, \mu_2, \ldots)$ is a filling of the
boxes of $\nu/\lambda$ with positive integers such that:
\begin{enumerate}
\item[(i)] entries weakly increase along each row (left to right),
\item[(ii)] entries strictly increase down each column,
\item[(iii)] the number of entries equal to $i$ is $\mu_i$, for each $i$.
\end{enumerate}
\end{definition}

\begin{definition}[Lattice word]
The \emph{reverse reading word} of a tableau $T$ is obtained by
reading the entries of $T$ from right to left in each row, starting
from the top row.  A word $w_1 w_2 \cdots w_N$ is a \emph{lattice
word} (or \emph{ballot sequence}) if, for every prefix $w_1 \cdots
w_j$ and every letter $i$, the number of occurrences of $i$ in the
prefix is at least the number of occurrences of $i+1$.
\end{definition}

\begin{theorem}[Littlewood-Richardson rule; see {\cite[Ch.~5]{fulton-1997}}, {\cite[Ch.~7]{stanley-1999}}]
\label{thm:lr-rule}
The coefficient $c^{\nu}_{\lambda,\mu}$ equals the number of
semistandard Young tableaux of skew shape $\nu/\lambda$, content
$\mu$, whose reverse reading word is a lattice word.  Such tableaux
are called \emph{LR tableaux}.
\end{theorem}

\begin{example}[The coefficient $c^{(3,2,1)}_{(2,1),\,(2,1)} = 2$]
\label{ex:lr-tableaux}
The skew shape $(3,2,1)/(2,1)$ consists of three boxes: position
$(1,3)$ in the first row, position $(2,2)$ in the second row, and
position $(3,1)$ in the third row.  We seek fillings with content
$(2,1)$ --- two $1$s and one $2$ --- satisfying the semistandard and
lattice word conditions.

The two LR tableaux are:
\[
\young(\hfil\hfil1,\hfil1,2)
\qquad\text{and}\qquad
\young(\hfil\hfil1,\hfil2,1)
\]
where $\hfil$ denotes a box belonging to $\lambda = (2,1)$ (shaded or
omitted in the skew diagram).  In the first tableau the reverse
reading word is $1,1,2$ and in the second it is $1,2,1$; both are
lattice words.  The filling $\young(\hfil\hfil2,\hfil1,1)$ with
reverse reading word $2,1,1$ is \emph{not} a lattice word (the prefix
$2$ has more $2$s than $1$s), so it is excluded.  Hence
$c^{(3,2,1)}_{(2,1),(2,1)} = 2$.
\end{example}

\subsubsection*{Knutson-Tao-Woodward puzzles}

An entirely different model for the same coefficient was introduced by
Knutson, Tao, and Woodward \cite{knutson-tao-woodward-2004}.

\begin{definition}[KTW puzzle]
Fix a positive integer $N$.  An \emph{$N$-puzzle} is a tiling of an
equilateral triangle of side length $N$ by the following unit pieces,
each carrying labels $0$ or $1$ on its edges:
\begin{enumerate}
\item A triangle of side $1$ with all three edges labeled $0$.
\item A triangle of side $1$ with all three edges labeled $1$.
\item A \emph{rhombus} (two unit triangles glued along a shared edge)
with edges labeled $0$ and $1$ in the pattern: the two edges adjacent
to one acute angle are labeled $1$, and the two edges adjacent to the
other acute angle are labeled $0$.
\end{enumerate}
The three rhombus orientations (at $0^\circ$, $60^\circ$, and
$120^\circ$) are all permitted.  In a valid tiling, adjacent pieces
must agree on shared edges.

The boundary of the triangle consists of $3N$ edges, which we read as
three binary strings of length~$N$: the \textbf{south} side (left to
right), the \textbf{northeast} side (bottom to top), and the
\textbf{northwest} side (top to bottom).  A partition $\lambda$
fitting inside a $k \times (N-k)$ box is encoded as a binary string
of length $N$ via the boundary-path encoding described in
Section~\ref{sec:classical-hopf} below.
\end{definition}

\begin{theorem}[{\cite{knutson-tao-woodward-2004}}]
\label{thm:ktw}
For partitions $\lambda, \mu, \nu$ fitting inside a $k \times (N-k)$
box, the LR coefficient $c^{\nu}_{\lambda,\mu}$ equals the number of
$N$-puzzles whose south boundary encodes $\nu$, northwest boundary
encodes $\lambda$, and northeast boundary encodes $\mu$.
\end{theorem}

Each valid tiling is a \emph{puzzle witness}: a self-contained
combinatorial proof that one unit of $c^{\nu}_{\lambda,\mu}$ exists.
For our running example, the two tableaux of
Example~\ref{ex:lr-tableaux} correspond to two puzzle tilings of a
triangle of side $N=5$.  We will see in Section~3 that each tiling is
a \emph{configuration} of an integrable lattice model --- a
frozen-in-place snapshot of particle worldlines threading through the
triangle.

\begin{remark}
Purbhoo \cite{purbhoo-2007} established a direct bijection between
KTW puzzles and LR tableaux via an intermediate object called a
\emph{mosaic}, making the equivalence of Theorems~\ref{thm:lr-rule}
and~\ref{thm:ktw} constructive rather than merely numerical.
\end{remark}

The multiplicity of models is the paper's motivating mystery.  Why
should rhombus tilings (a geometric object), lattice-word conditions
on skew tableaux (a combinatorial object), and tensor product
multiplicities (a representation-theoretic object) all count the same
thing?  One answer --- the one this paper develops --- is that all
three are shadows of a single algebraic structure: an exactly solvable
lattice model whose transfer matrices commute by the Yang-Baxter
equation.

\subsection{The Hopf algebra perspective}
\label{sec:classical-hopf}

The three definitions of Section~\ref{sec:classical-defn} present LR
coefficients as \emph{structure constants} of a product.  But the ring
of symmetric functions carries more structure: it is a
\emph{self-dual graded connected Hopf algebra}, and the LR
coefficients appear equally in the coproduct.  This dual perspective
is less standard in the combinatorics literature, but it is the
natural home for transfer operators and the gateway to integrability.

\subsubsection*{The coproduct}

The ring $\Lambda$ of symmetric functions carries a coproduct
\[
  \Delta \colon \Lambda \to \Lambda \otimes \Lambda
\]
defined on the power-sum generators by $\Delta(p_k) = p_k \otimes 1 +
1 \otimes p_k$ and extended as an algebra homomorphism.  On the Schur
basis, the coproduct takes the form
\begin{equation}
\label{eq:coproduct}
  \Delta(s_\lambda) = \sum_{\mu,\nu} c^{\lambda}_{\mu,\nu}\, s_\mu
  \otimes s_\nu,
\end{equation}
where the same LR coefficients appear, now with $\lambda$ upstairs
and the summation over all pairs $(\mu, \nu)$ with $|\mu| + |\nu| =
|\lambda|$.  The coproduct encodes \emph{all} LR coefficients for a
given $\lambda$ at once: it ``unfolds'' the partition like a flower
opening, each summand a valid way to split the whole into two
interlocking halves.

Self-duality means that the structure constants of the product and
coproduct coincide: the coefficient of $s_\mu \otimes s_\nu$ in
$\Delta(s_\lambda)$ equals $c^\lambda_{\mu,\nu}$, the same integer
that appears as the coefficient of $s_\lambda$ in $s_\mu \cdot s_\nu$.
This is the Hopf-algebraic shadow of the fact that the LR
coefficients are ``the same number'' regardless of which definition we
use.

\subsubsection*{Binary string encoding}

To make the coproduct computable by transfer operators, we need a
finite-dimensional encoding of partitions.

\begin{definition}[Binary string encoding]
\label{def:binary-encoding}
Fix a bounding box of height $k$ (number of parts) and width $m$
(largest part).  A partition $\lambda$ with at most $k$ parts, each at
most $m$, is encoded as a binary string $b(\lambda) \in \{0,1\}^{k+m}$
by tracing the northeast boundary of the Young diagram of $\lambda$
inside the $k \times m$ box, from the southwest corner to the
northeast corner:
\begin{itemize}
\item a vertical step (north) is recorded as $1$,
\item a horizontal step (east) is recorded as $0$.
\end{itemize}
The string has exactly $k$ ones and $m$ zeros.  For instance, in a
$3 \times 4$ box, the partition $(3,2,1)$ traces the boundary path
$\rightarrow \uparrow \rightarrow \uparrow \rightarrow \uparrow
\rightarrow$, giving $b(3,2,1) = 0101010$.
\end{definition}

This encoding maps the lattice of partitions inside the $k \times m$
box to a subset of the Boolean lattice $\{0,1\}^{k+m}$.  It is also
the encoding used for the boundary conditions of KTW puzzles: the
binary strings on the three sides of the puzzle triangle are precisely
these boundary-path encodings of $\lambda$, $\mu$, and $\nu$.

\subsubsection*{Transfer operators}

With partitions encoded as binary strings, we can define operators
that act on formal linear combinations of partitions by stripping away
horizontal or vertical strips --- one row of boxes at a time,
interlocking with the previous state like teeth of a zipper.

\begin{definition}[Free-fermionic transfer operators]
Let $\lambda$ be a partition inside a $k \times m$ box.
\begin{enumerate}
\item $T_{\mathrm{free}}(u)\,|\lambda\rangle = \displaystyle\sum_{\substack{\mu
\subseteq \lambda \\ \lambda/\mu \text{ h-strip}}} u^{|\lambda/\mu|}
\,|\mu\rangle$, where the sum is over all $\mu$ such that
$\lambda/\mu$ is a \emph{horizontal strip} (at most one box per
column).  In terms of binary strings, $T_{\mathrm{free}}$ slides each
$0$ (red particle) rightward past one or more $1$s (green particles);
each crossing contributes a factor of $u$.

\item $\bar{T}_{\mathrm{free}}(u)\,|\lambda\rangle =
\displaystyle\sum_{\substack{\mu \subseteq \lambda \\ \lambda/\mu
\text{ v-strip}}} u^{|\lambda/\mu|}\, |\mu\rangle$, where the sum is
over all $\mu$ such that $\lambda/\mu$ is a \emph{vertical strip} (at
most one box per row).  In the binary string picture, each $1$ (green
particle) moves leftward by at most one step past a $0$, each move
contributing a factor of $u$.
\end{enumerate}
\end{definition}

The key identity connecting these operators to symmetric functions is:
\begin{equation}
\label{eq:schur-transfer}
  s_{\lambda/\mu}(u_1, \ldots, u_n)
  = \langle \mu \,|\, T_{\mathrm{free}}(u_1) \cdots
    T_{\mathrm{free}}(u_n)\, |\lambda\rangle,
\end{equation}
where $\langle\mu|$ denotes the linear functional picking out the
coefficient of $|\mu\rangle$.  Setting $\mu = \varnothing$ recovers
the Schur function $s_\lambda(u_1, \ldots, u_n)$.  The analogous
formula with $\bar{T}_{\mathrm{free}}$ computes the \emph{conjugate}
skew Schur function $s_{\lambda'/\mu'}$.

The coproduct~\eqref{eq:coproduct} can now be computed by composing
transfer operators.  Write $T_{\mathrm{free}}$ applied $n$ times (with
variables $u_1, \ldots, u_n$) to expand
$\Delta(s_\lambda)$: one obtains $s_{\lambda/\mu}$ as a polynomial in
the $u_i$, and then decomposes it into Schur functions to read off the
LR coefficients.

\subsubsection*{A computational subtlety}

In practice, extracting individual LR coefficients from the transfer
operator output requires care.  The na\"ive approach --- evaluate the
Schur function $s_\nu(1, \ldots, 1)$ at various numbers of variables
using the hook-content formula, then solve a linear system --- fails
because the hook-content values are \emph{linearly dependent} across
partitions of the same size.  For example,
\[
  \dim_n(3,1) - 3\,\dim_n(2,2) + \dim_n(2,1,1) = 0
  \quad \text{for all } n,
\]
where $\dim_n(\nu) = s_\nu(1^n)$ is the dimension of the $GL_n$
representation $V_\nu$.  So the principal specialization cannot
separate all Schur functions of a given degree.

The fix, discovered through computational experiment
\cite{zinn-justin-2008}, is to use the \emph{Kostka matrix}.  Recall
that the Kostka number $K_{\nu,\rho}$ counts semistandard Young
tableaux of shape $\nu$ and content $\rho$.  Ordered by dominance, the
Kostka matrix is \emph{unitriangular}: $K_{\nu,\nu} = 1$ and
$K_{\nu,\rho} = 0$ unless $\nu \trianglerighteq \rho$ in dominance
order.  In particular, the Kostka matrix is always invertible, and its
inverse can be computed by forward substitution.

One computes the \emph{skew Kostka numbers}
$\widetilde{K}_{\lambda/\mu,\,\rho}$ (counting SSYT of skew shape
$\lambda/\mu$ and content $\rho$) via chains of $T_{\mathrm{free}}$
applications --- each chain of strip removals of prescribed sizes
traces a single SSYT.  Then, inverting the Kostka matrix recovers the
individual LR coefficients $c^\lambda_{\mu,\nu}$ from the skew Kostka
numbers via
\[
  \widetilde{K}_{\lambda/\mu,\,\rho}
  = \sum_\nu c^\lambda_{\mu,\nu}\, K_{\nu,\rho}.
\]
The unitriangularity of $K$ makes this inversion exact over the
integers, with no linear dependencies to contend with.

\begin{example}[Coproduct of $s_{(3,2,1)}$]
\label{ex:coproduct-321}
Applying the transfer operator machinery to $\lambda = (3,2,1)$
yields the coproduct
\[
  \Delta(s_{(3,2,1)}) = s_{(3,2,1)} \otimes s_\varnothing
  + s_{(3,1)} \otimes s_{(1,1)}
  + s_{(2,2)} \otimes s_{(1,1)}
  + s_{(3,2)} \otimes s_{(1)}
  + \cdots
  + 2\, s_{(2,1)} \otimes s_{(2,1)}
  + \cdots
  + s_\varnothing \otimes s_{(3,2,1)},
\]
with $25$ nonzero terms in total.  The coefficient $2$ at
$s_{(2,1)} \otimes s_{(2,1)}$ is our running example.
\end{example}

\subsubsection*{From transfer operators to lattice models}

We close with the observation that motivates the rest of this paper.
The operators $T_{\mathrm{free}}$ and $\bar{T}_{\mathrm{free}}$ are
\emph{row transfer matrices} of a lattice model.  Each application of
$T_{\mathrm{free}}$ processes one row of the lattice: the input state
(a binary string encoding a partition) enters from below, the operator
produces all valid output states above, weighted by powers of the
spectral parameter $u$.  Composing $n$ transfer matrices builds the
lattice row by row, and the matrix element
$\langle\varnothing|T_{\mathrm{free}}(u_1) \cdots
T_{\mathrm{free}}(u_n)|\lambda\rangle$ is a partition function: a sum
over all configurations of the lattice with the prescribed boundary
conditions.

In Section~3 we will see that this is not merely an analogy.  The
puzzle pieces of Knutson-Tao-Woodward are precisely the
\emph{R-matrix weights} of a free-fermionic vertex model, and the row
transfer matrices commute because the R-matrix satisfies the
Yang-Baxter equation.  The multiplicity of combinatorial models ---
puzzles, tableaux, mosaics --- is then no mystery: they are different
ways of summing over the configurations of a single integrable system.
